\section{Scriptural Date and Location}\label{sec:date-location}
% Generated. Do not edit.
% Deterministic projection of research/chronology.toml.
% schema: 4
% document: liturgy/roman-rite/1962/propers/temporal/60-twentieth-after-pentecost
% generated-by: tools/tpt proper-chronology annotations
\newcommand{\chronologyannotationclaim}[6]{#6}
\newcommand{\chronologyannotationcomparisonclaim}[8]{#8}
\newcommand{\chronologyannotationanchorclaim}[9]{#9}
\newcommand{\chronologyannotationanchorgroup}[8]{#8}
\newcommand{\chronologyannotationreach}[3]{}
\newcommand{\chronologyannotationgroup}[3]{#3}
\newcommand{\chronologyannotationcomparisongroup}[4]{#4}
\newcommand{\chronologyannotation}[1]{%
  \ifcsname triptychchronologyannotation@#1\endcsname
    \csname triptychchronologyannotation@#1\endcsname
  \else
    \PackageError{triptych}{No chronology annotation for #1}{}%
  \fi}
% comparison-dependency: src/gpt/liturgy/roman-rite/1962/propers/temporal/60-twentieth-after-pentecost/research/chronology-profile-comparisons.toml sha256=11a4bde1a39800640daad7a03563050ab63f20b33c038c8773a288bf407b2b9d
% introit: Introit
\expandafter\def\csname triptychchronologyannotation@introit\endcsname{%
  \chronologyannotationgroup{historical-setting}{nonuniform}{\textbf{Historical setting} -- No single occasion date applies across every cited locus.}%
  \space \chronologyannotationgroup{composition}{nonuniform}{\textbf{Composition} -- No single date applies across every cited locus.}%
}
% epistle: Epistle
\expandafter\def\csname triptychchronologyannotation@epistle\endcsname{%
  \chronologyannotationgroup{composition}{disputed}{\textbf{Composition} -- disputed: \chronologyannotationreach{Eph.5.15}{Eph}{true}\chronologyannotationreach{Eph.5.16}{Eph}{true}\chronologyannotationreach{Eph.5.17}{Eph}{true}\chronologyannotationreach{Eph.5.18}{Eph}{true}\chronologyannotationreach{Eph.5.19}{Eph}{true}\chronologyannotationreach{Eph.5.20}{Eph}{true}\chronologyannotationreach{Eph.5.21}{Eph}{true}\chronologyannotationclaim{composition.epistle-to-the-ephesians}{composition}{catholic-traditional-v1}{disputed}{a period between 58 and 63}{A.D. 58--63}; \chronologyannotationreach{Eph.5.15}{Eph}{true}\chronologyannotationreach{Eph.5.16}{Eph}{true}\chronologyannotationreach{Eph.5.17}{Eph}{true}\chronologyannotationreach{Eph.5.18}{Eph}{true}\chronologyannotationreach{Eph.5.19}{Eph}{true}\chronologyannotationreach{Eph.5.20}{Eph}{true}\chronologyannotationreach{Eph.5.21}{Eph}{true}\chronologyannotationclaim{composition.epistle-to-the-ephesians}{composition}{catholic-traditional-v1}{disputed}{(Philemon; Colossians; Ephesians; Philippians), 61}{A.D. 61}.}%
  \space \chronologyannotationcomparisongroup{catholic-critical-v1}{composition}{composition-only}{\textbf{Composition (Catholic critical chronology, for comparison)} -- disputed: \chronologyannotationreach{Eph.5.15}{Eph}{true}\chronologyannotationreach{Eph.5.16}{Eph}{true}\chronologyannotationreach{Eph.5.17}{Eph}{true}\chronologyannotationreach{Eph.5.18}{Eph}{true}\chronologyannotationreach{Eph.5.19}{Eph}{true}\chronologyannotationreach{Eph.5.20}{Eph}{true}\chronologyannotationreach{Eph.5.21}{Eph}{true}\chronologyannotationcomparisonclaim{critical.ephesians-later-disciple}{composition}{catholic-critical-v1}{catholic-critical-v1}{disputed}{passage.united-states-conference-of-catholic-bishops.new-american-bible-revised-edition.english-usccb-web-2026-09-28.ephesians-introduction}{around A.D. 80--100}{Ephesians under the later-disciple hypothesis (approximate), c. A.D. 80--100}.}%
}
% gradual: Gradual
\expandafter\def\csname triptychchronologyannotation@gradual\endcsname{%
  \chronologyannotationgroup{traditional-attribution}{disputed}{\textbf{Traditional attribution} -- disputed: \chronologyannotationreach{Ps.144.15}{Ps.107 Ps.144}{true}\chronologyannotationreach{Ps.144.16}{Ps.107 Ps.144}{true}\chronologyannotationclaim{israel.monarchy.david-traditional-era}{traditional-attribution}{catholic-traditional-v1}{disputed}{reigned from 1055 to 1015 B.C.}{David (reign in the usual chronology), B.C. 1055--1015}.}%
  \space \chronologyannotationgroup{historical-setting}{research-pending}{\textbf{Historical setting} -- Occasion date unresolved.}%
  \space \chronologyannotationgroup{composition}{preferred}{\textbf{Composition}: \chronologyannotationreach{Ps.144.15}{Ps}{true}\chronologyannotationreach{Ps.144.16}{Ps}{true}\chronologyannotationclaim{critical.psalms.latest-composition-boundary}{composition}{catholic-critical-v1}{preferred}{before the Maccabean period, around 165 B.C.; no individual psalm can be dated securely}{Before c. 165 B.C.}}%
}
% alleluia: Alleluia
\expandafter\def\csname triptychchronologyannotation@alleluia\endcsname{%
  \chronologyannotationgroup{traditional-attribution}{disputed}{\textbf{Traditional attribution} -- disputed: \chronologyannotationreach{Ps.107.2}{Ps.107 Ps.144}{true}\chronologyannotationclaim{israel.monarchy.david-traditional-era}{traditional-attribution}{catholic-traditional-v1}{disputed}{reigned from 1055 to 1015 B.C.}{David (reign in the usual chronology), B.C. 1055--1015}.}%
  \space \chronologyannotationgroup{historical-setting}{research-pending}{\textbf{Historical setting} -- Occasion date unresolved.}%
  \space \chronologyannotationgroup{composition}{preferred}{\textbf{Composition}: \chronologyannotationreach{Ps.107.2}{Ps}{true}\chronologyannotationclaim{critical.psalms.latest-composition-boundary}{composition}{catholic-critical-v1}{preferred}{before the Maccabean period, around 165 B.C.; no individual psalm can be dated securely}{Before c. 165 B.C.}}%
}
% gospel: Gospel
\expandafter\def\csname triptychchronologyannotation@gospel\endcsname{%
  \chronologyannotationgroup{narrated-event}{preferred}{\textbf{Event}: \chronologyannotationreach{John.4.46}{John.4.46-53}{false}\chronologyannotationreach{John.4.47}{John.4.46-53}{false}\chronologyannotationreach{John.4.48}{John.4.46-53}{false}\chronologyannotationreach{John.4.49}{John.4.46-53}{false}\chronologyannotationreach{John.4.50}{John.4.46-53}{false}\chronologyannotationreach{John.4.51}{John.4.46-53}{false}\chronologyannotationreach{John.4.52}{John.4.46-53}{false}\chronologyannotationreach{John.4.53}{John.4.46-53}{false}\chronologyannotationclaim{life-of-christ.healing-of-the-rulers-son-at-cana}{narrated-event}{catholic-traditional-v1}{preferred}{Passover, A.U.C. 779 - about Pentecost, 780}{The healing of the ruler's son at Cana, A.D. 26--27 (derived)}.}%
  \space \chronologyannotationgroup{composition}{disputed}{\textbf{Composition} -- disputed: \chronologyannotationreach{John.4.46}{John}{true}\chronologyannotationreach{John.4.47}{John}{true}\chronologyannotationreach{John.4.48}{John}{true}\chronologyannotationreach{John.4.49}{John}{true}\chronologyannotationreach{John.4.50}{John}{true}\chronologyannotationreach{John.4.51}{John}{true}\chronologyannotationreach{John.4.52}{John}{true}\chronologyannotationreach{John.4.53}{John}{true}\chronologyannotationclaim{composition.gospel-of-john}{composition}{catholic-traditional-v1}{disputed}{from the year 90 to 100 (approximately)}{c. A.D. 90--100}; \chronologyannotationreach{John.4.46}{John}{true}\chronologyannotationreach{John.4.47}{John}{true}\chronologyannotationreach{John.4.48}{John}{true}\chronologyannotationreach{John.4.49}{John}{true}\chronologyannotationreach{John.4.50}{John}{true}\chronologyannotationreach{John.4.51}{John}{true}\chronologyannotationreach{John.4.52}{John}{true}\chronologyannotationreach{John.4.53}{John}{true}\chronologyannotationclaim{composition.gospel-of-john}{composition}{catholic-traditional-v1}{disputed}{96 or one of the succeeding years}{A.D. 96--100}.}%
}
% offertory: Offertory
\expandafter\def\csname triptychchronologyannotation@offertory\endcsname{%
  \chronologyannotationgroup{historical-setting}{preferred}{\textbf{Historical setting}: \chronologyannotationreach{Ps.136.1}{Ps.136}{true}\chronologyannotationclaim{israel.exile.psalm-136-captive-lament}{historical-setting}{catholic-traditional-v1}{preferred}{For the time of Jeremias, and the captivity of Babylon}{The Babylonian captives' lament, after Jerusalem's destruction under Sedecias (derived)}.}%
  \space \chronologyannotationgroup{historical-setting}{preferred}{\textbf{Historical setting}: \chronologyannotationreach{Ps.136.1}{Ps.136}{true}\chronologyannotationclaim{israel.exile.seventy-years}{historical-setting}{catholic-traditional-v1}{preferred}{seventy years}{The seventy years of servitude to Babylon, ``seventy years'' (duration)}.}%
  \space \chronologyannotationgroup{composition}{preferred}{\textbf{Composition}: \chronologyannotationreach{Ps.136.1}{Ps}{true}\chronologyannotationclaim{critical.psalms.latest-composition-boundary}{composition}{catholic-critical-v1}{preferred}{before the Maccabean period, around 165 B.C.; no individual psalm can be dated securely}{Before c. 165 B.C.}}%
  \space \chronologyannotationanchorgroup{israel.exile.psalm-136-captive-lament}{historical-setting}{catholic-traditional-v1}{For the time of Jeremias, and the captivity of Babylon}{israel.exile.third-captivity}{after}{context}{\textbf{Historical background: The third captivity of Juda: the destruction of Jerusalem under Sedecias}: Preferred: \chronologyannotationanchorclaim{israel.exile.third-captivity}{0}{catholic-traditional-v1}{preferred}{passage.george-leo-haydock.douay-rheims-with-haydock-commentary.2014-loreto-feeney-memorial.psalm-70-captivities-usher-chronology}{reported-excluded}{A.M. 3416}{Haydock reporting Ussher; A.M. is not converted}{A.M. 3416 (Haydock reporting Ussher; A.M. is not converted)}; alternatives: \chronologyannotationanchorclaim{israel.exile.third-captivity}{1}{catholic-traditional-v1}{alternate}{artifact.catholic-encyclopedia.volume-3.new-york-1908.newadvent-03315a-9cb437e2,artifact.catholic-encyclopedia.volume-14.new-york-1912.newadvent-14499a-82b12b6e}{traditional-catholic}{586 B.C.}{Reid and Meistermann, Catholic Encyclopedia}{B.C. 586 (Reid and Meistermann, Catholic Encyclopedia)}, \chronologyannotationanchorclaim{israel.exile.third-captivity}{2}{catholic-traditional-v1}{alternate}{artifact.catholic-encyclopedia.volume-8.new-york-1910.newadvent-08652a-189e1a8f}{traditional-catholic}{the fall of Jerusalem (587 B.C.)}{Schets, Catholic Encyclopedia}{B.C. 587 (Schets, Catholic Encyclopedia)}, \chronologyannotationanchorclaim{israel.exile.third-captivity}{3}{catholic-traditional-v1}{alternate}{artifact.catholic-encyclopedia.volume-8.new-york-1910.newadvent-08654a-645bba6c}{reported-traditional}{B.C. 588}{Petavius table reported by Sloet; retained transcription checked, printed-page verification still pending}{B.C. 588 (Petavius table reported by Sloet; retained transcription checked, printed-page verification still pending)}. The represented event is after this historical event. These dates belong to that historical event, not to the passage or its composition. Both events could occur in the same year.}%
}
% communion: Communion
\expandafter\def\csname triptychchronologyannotation@communion\endcsname{%
  \chronologyannotationgroup{traditional-attribution}{disputed}{\textbf{Traditional attribution} -- disputed: \chronologyannotationreach{Ps.118.49}{Ps.118}{true}\chronologyannotationreach{Ps.118.50}{Ps.118}{true}\chronologyannotationclaim{israel.monarchy.david-traditional-era}{traditional-attribution}{catholic-traditional-v1}{disputed}{reigned from 1055 to 1015 B.C.}{David (reign in the usual chronology), B.C. 1055--1015}.}%
  \space \chronologyannotationgroup{historical-setting}{research-pending}{\textbf{Historical setting} -- Occasion date unresolved.}%
  \space \chronologyannotationgroup{composition}{preferred}{\textbf{Composition}: \chronologyannotationreach{Ps.118.49}{Ps}{true}\chronologyannotationreach{Ps.118.50}{Ps}{true}\chronologyannotationclaim{critical.psalms.latest-composition-boundary}{composition}{catholic-critical-v1}{preferred}{before the Maccabean period, around 165 B.C.; no individual psalm can be dated securely}{Before c. 165 B.C.}}%
}


The passages appear in Catholic canonical and verse order, with Vulgate
psalm numbering first. The dates distinguish attribution, historical
setting, and composition. A received author's reference era does not
date a poem's particular occasion. The Psalter-wide composition boundary
comes from the USCCB introduction's broader account; it is not an
individual estimate for each appointed psalm.

\begin{dossiertable}
Alleluia & Ps.~107:2 (108) & Davidic title; Israel's praise and threatened dependence & \chronodate{alleluia}{\chronologyannotation{alleluia}}\\
\cmidrule(lr){1-4}
\dossierprose{The inherited title names David. Corbett's customary regnal chronology supplies the reference era, with its disputed status; neither it nor the psalm's reuse of earlier material locates a composition site or biographical occasion. Augustine, Ps.~107.1--2, recognizes the reused psalm passages. The USCCB Psalms introduction denies secure dating of individual poems.}
\midrule
Introit verse & Ps.~118:1 (119) & Israel's prayer under the law; received Davidic attribution & \chronodate{introit}{\chronologyannotation{introit}}\\
\cmidrule(lr){1-4}
\dossierprose{This appointment also contains the Daniel antiphon, whose separate setting appears below; hence no date governs the Introit as a whole. Ps~118:1: B.C.~1055--1015 is David's reign as a traditional-attribution reference; Ps~118:1: before c.~165~B.C. is the broad Psalter composition boundary. Haydock's note at 118:1 reports Davidic attribution, instruction of Solomon, David's persecution by Saul, and a proposal associating Daniel with the captives' consolation. He favors David over conjectures of another author. These alternatives locate no single writing occasion or site.}
\midrule
Communion & Ps.~118:49--50 (119) & The servant's hope amid humiliation; the same poem and attribution & \chronodate{communion}{\chronologyannotation{communion}}\\
\cmidrule(lr){1-4}
\dossierprose{The promise and humiliation belong to the poem's extended prayer for faithful life. Haydock's received authorship and occasion alternatives apply here as to its opening verse. The displayed regnal span remains a reference for the traditional attribution, not a date for this humiliation. The composition boundary applies to the Psalter's broad formation, not an identified episode in the speaker's life.}
\midrule
Offertory & Ps.~136:1 (137) & Babylonian captivity; Jerusalem remembered by displaced Israelites & \chronodate{offertory}{\chronologyannotation{offertory}}\\
\cmidrule(lr){1-4}
\dossierprose{Haydock, Ps.~136, preserves variant or absent headings and proposals involving captivity, return, and prophetic composition. The represented lament and the destruction it remembers are not the same event; neither by itself identifies the poem's writing date or site. Jeremiah~25 supplies the servitude's duration, not an absolute date for the lament. The background alternatives above date Jerusalem's destruction. Haydock reports Ussher's era; Sloet reports Petavius's table, with the displayed verification limit. Reid and Meistermann differ from Schets. The alternatives remain separate, not a merged interval or a newly calculated lower bound for this verse.}
\midrule
Gradual & Ps.~144:15--16 (145) & Davidic praise; creatures dependent on the divine king & \chronodate{gradual}{\chronologyannotation{gradual}}\\
\cmidrule(lr){1-4}
\dossierprose{The psalm's Davidic title supplies traditional attribution. Its universal nourishment and kingdom language name neither a particular incident nor a certain place of composition. Corbett supplies the customary regnal reference; the USCCB introduction supplies only the broader Psalter boundary. Christian Eucharistic reception belongs to the passage's later liturgical life.}
\midrule
Introit antiphon & Dan.~3:29,31,35 & Azarias and his companions in Babylon's furnace & Prayer's composition: most likely long after the Exile (Gigot); represented setting: servitude to Babylon.\\
\cmidrule(lr){1-4}
\dossierevent{Narrated event: Daniel~3:23--25 places Azarias's prayer within the furnace after refusal of the king's image. The servitude's ``seventy years'' is a duration supplied by Jeremiah~25, not the furnace episode's length or a calendar year.}
\cmidrule(lr){1-4}
\dossierprose{Gigot, ``Book of Daniel,'' concluding section on the deuterocanonical portions, distinguishes the Prayer and Song from whole-book composition theories. He reports widespread early attribution but favors an unknown Jewish author long after the Exile for this portion. That relative proposal concerns writing, not the narrated event. The prayer speaks within Israel's displaced worship and covenant memory; no more precise author, composition site, or life stage is established. The Introit adapts and reorders its verses.}
\midrule
Gospel & John~4:46b--53 & Jesus at Cana; child at Capharnaum, both in Galilee; traditional Johannine writing associated with Ephesus & \chronodate{gospel}{\chronologyannotation{gospel}}\\
\cmidrule(lr){1-4}
\dossierevent{Narrated event: Cana is named in the unappointed first part of John~4:46. The father seeks Jesus there; his child is at Capharnaum. Maas's encompassing journey supplies the derived event range. It is not an absolute date printed beside the healing, and the matching seventh hour does not supply a modern clock-time.}
\cmidrule(lr){1-4}
\dossierprose{Fonck and Durand give the displayed traditional composition alternatives for the Gospel, associated with John's later life and an audience beyond the original event. Composition in the Johannine tradition and the event in Galilee are separate horizons. The healing's report and the household's belief belong to the narrative; the work's publication is not dated by their hour.}
\midrule
Epistle & Eph.~5:15--21 & Traditional Pauline captivity; Rome favored over Caesarea in Ladeuze; recipients associated with Ephesus and Asia Minor & \chronodate{epistle}{\chronologyannotation{epistle}}\\
\cmidrule(lr){1-4}
\dossierprose{Ladeuze's captivity range and Prat's narrower chronology are retained as alternatives. The USCCB Ephesians introduction distinguishes direct Pauline authorship, composition through a secretary, and a later disciple's development of Pauline teaching; style, vocabulary, ecclesiology, and comparison with Colossians inform the debate. The critical comparison's approximate interval dates only the later-disciple proposal, not the secretarial hypothesis. The impersonal address and omitted place-name in early witnesses leave destination debated. Proposed Asian circulation does not identify the writer's location.}
\end{dossiertable}

These chronology proposals concern distinguishable claims and witnesses.
The source locators in the references identify the accounts behind them;
the date of the present liturgical occurrence is a separate fact.
